\section{Problem Setup and Benchmark}
\label{sec:benchmark}

\subsection{Setting and Metrics}
\label{sec:setup-metrics}
A detector produces a window-level anomaly score $\st$ from a multivariate window $\xt$ (window length 256 samples, stride 128, native sampling rates throughout). Training uses \emph{source-normal} windows only, pooled across source domains. At deployment, the target stream is unlabeled; we may use an initial segment as an unlabeled calibration stream but never its labels. A calibration rule turns scores into thresholds $\qt$ and alarms $\mathbb{1}[\st > \qt]$.

We report the following quantities. \emph{Normal-window coverage} $\cover = 1-\FPR$ is the fraction of evaluation windows with label ``normal'' that are not alarmed, with the requirement $\cover \approx 1-\alphaNom$ at nominal level $\alphaNom$; unless stated otherwise, coverage is the calibration target and all calibration tables report it. The \emph{coverage gap} is $\gap = |\cover-(1-\alphaNom)|$. We also report the \emph{alarm rate} over the whole stream, because an online rule that regulates alarms over a stream containing anomalies cannot in general also control the normal-window FPR: the two coincide only when anomalies are rare relative to $\alphaNom$. \emph{Interval inflation} is $\overline{\qt}/\qsrc$, the mean deployed threshold relative to the source-calibrated quantile; it measures the transfer cost of keeping coverage. Detection-side metrics are range-based $F_1$, AUROC/AUPRC, false alarms per window (equivalently per 1000 windows), and detection delay in windows after an anomaly run begins. All calibration results are averaged over three training seeds; we add bootstrap 95\% confidence intervals on the coverage gap computed by resampling evaluation windows.

\subsection{Six Folds over Four Domain Families}
\label{sec:folds}
We construct six leave-one-group-out folds (each fold holds out one dataset group; four domain families are covered). Table~\ref{tab:datasets} lists targets, native rates, evaluation set sizes, anomaly fractions, and label provenance; sources for each fold are all other groups' normal windows. Two guardrails apply throughout: folds whose ranking quality is below chance are excluded from detection-quality claims (diagnosed in Section~\ref{sec:decoupling}), and the synthetic-injection fold (F3) is flagged in every table and figure caption. Two dataset-engineering decisions follow from pilot experiments and are fixed throughout: (i) windows are standardized with per-dataset global statistics---per-window standardization removes the amplitude component of degradation and we observed it destroy ranking on run-to-failure data; (ii) no resampling to a common rate is performed, so spectral differences between domains are preserved.

\begin{table}[t]
\caption{Folds, datasets, native rates, and evaluation sets (windows of length 256, stride 128). Anomaly fractions vary by protocol and matter for interpretering alarm-budget effects.}
\label{tab:datasets}
\centering
\small
\begin{tabular}{llrrrr}
\toprule
Fold & Target (rate) & Rate & Normal & Anom. & Label source \\
\midrule
F1 & PHM16-CMP (1\,Hz)        & 1\,Hz          & 1912 & 251   & removal-rate $q_{10}$ drift \\
F2 & CARE wind (10\,min)      & 10\,min        & 4714 & 557   & status labels \\
F3 & PVDAQ solar ($\sim$10--15\,min) & $\sim$12\,min & 858 & 214 & synthetic injection \\
F4 & CWRU+Paderborn           & 12/128\,kHz    & 10980 & 7556 & seeded / last-20\% \\
F5 & MFPT                     & 97.7\,kHz      & 5072 & 13661 & segment labels \\
F6 & XJTU-SY                  & 25.6\,kHz      & 12806 & 3993 & last-20\% run-to-failure \\
\bottomrule
\end{tabular}
\end{table}

\subsection{Labeling and Injection Protocols}
\label{sec:protocols}
Pseudo-labeling is unavoidable in parts of this benchmark, and each protocol is a documented limitation. PHM16-CMP \cite{phm2016} has no fault labels; we mark a wafer-stage run anomalous when its removal rate falls in the lowest per-stage decile ($q_{10}$), giving an 11.6\% anomaly fraction. CARE wind \cite{care2024} carries status-based anomaly labels that we use directly. PVDAQ \cite{pvdaq} contains normal-operation data only; the F3 fold evaluates a synthetic injection protocol---$+6\sigma$ spikes, $+4\sigma$ level shifts, and segment dropouts applied to 20\% of evaluation windows with a fixed seed---and is reported as synthetic-only. CWRU supplies seeded-fault files (label normal/fault); Paderborn \cite{paderborn} contributes run-to-failure records with the last 20\% of measurement indices labeled anomalous. MFPT segments carry per-segment labels. XJTU-SY \cite{wang2020} follows the last-20\% run-to-failure convention. All folds share the same encoder and calibration code; dataset sizes are equalized by deterministic window subsampling (seed fixed) where needed.

\subsection{Encoder and Training}
\label{sec:encoder}
Scores come from a plain PatchTST-style reconstruction model \cite{patchtst2023}: patch length 64, stride 32, $d_{\mathrm{model}}=64$, 4 heads, 2 encoder layers, band-conditioned patch embeddings; the per-window score is the mean squared reconstruction error. Models train for 4000 steps with batch size 64 (AdamW, learning rate $5\times10^{-4}$, cosine schedule, AMP) on pooled source-normal windows; target scores are produced zero-shot. No test-time adaptation, gating, or invariance regularizers are used in the main runs: earlier ablations showed their contributions were marginal ($\leq 2.2\%$ range-$F_1$ changes across all four Stage-0 folds) and they are reported as negative controls in Section~\ref{sec:neg-controls}.
